\chapter{Cybersecurity and Data Protection Risk Assessment}
\label{chap:Cybersecurity and Data Protection Risk Assessment}

%---------------------------- SECTION ----------------------------
\section{The Risk That Touches Everything}
\paragraph{There is a category of risk that, within the space of a single professional generation, has moved from the periphery of organisational concern to its very centre. A risk that affects every organisation regardless of size, sector, or geography. A risk that evolves faster than the regulatory frameworks designed to address it, faster than the controls designed to manage it, and faster than the understanding of most boards and senior management teams responsible for overseeing it.}
\paragraph{That risk is cybersecurity — and its inseparable companion in the compliance world, data protection.}
\paragraph{Twenty years ago, cybersecurity was primarily a technical concern — the domain of IT departments managing firewalls and antivirus software. Data protection was a compliance obligation of modest practical significance for most organisations. Neither was typically on the agenda of boards, regulators, or senior compliance professionals in the way that credit risk, market risk, or operational risk might be.}
\paragraph{Today, the picture could not be more different. Cyber attacks are among the most significant operational risks facing organisations across every sector. Data breaches generate regulatory investigations, substantial fines, class action litigation, and reputational damage that can take years to repair. Ransomware attacks have paralysed hospitals, disrupted critical national infrastructure, and brought major corporations to a standstill. And the regulatory frameworks around both cybersecurity and data protection have expanded dramatically — imposing specific, enforceable, and increasingly detailed obligations on organisations to assess, manage, and report their cyber and data risks.}
\paragraph{For compliance and legal professionals, this is not someone else's territory. It is one of the most important and most actively developing areas of the compliance landscape — one that requires genuine engagement, genuine understanding, and a risk assessment framework that is both technically credible and organisationally embedded.}
\paragraph{This chapter provides that framework.}

%---------------------------- SECTION ----------------------------
\section{Understanding the Threat Landscape}
\paragraph{Effective cybersecurity risk assessment starts with a clear understanding of the threat landscape — the range of threats that could exploit vulnerabilities in the organisation's systems, processes, and people to cause harm. Unlike most other risk categories, the cyber threat landscape is both externally driven and continuously evolving — threat actors are adaptive, innovative, and motivated in ways that make the risk landscape genuinely dynamic.}

\subsection{Threat Actors — Who Is Behind the Risk?}
\paragraph{Understanding who poses a cyber threat — and what motivates them — is a fundamental input into cybersecurity risk assessment. Different threat actors have different capabilities, different objectives, and different methods — and the controls that are most effective against one type of threat may be less relevant to another.}
\paragraph{\textbf{Financially motivated criminal organisations} — organised criminal groups targeting organisations for financial gain represent the most prevalent category of cyber threat for most businesses. Ransomware — malicious software that encrypts an organisation's data and demands payment for its release — has become the most significant cyber threat facing organisations across every sector. Ransomware attacks have affected hospitals, schools, law firms, financial institutions, manufacturers, and government bodies — and the financial, operational, and reputational consequences have in many cases been severe and prolonged.}
\paragraph{\textbf{Nation-state actors} — state-sponsored cyber operations represent a qualitatively different category of threat — characterised by sophisticated capabilities, significant resources, long-term persistence, and strategic rather than purely financial objectives. Nation-state cyber operations may target critical national infrastructure, seek to steal intellectual property or sensitive government information, conduct espionage against specific organisations or individuals, or — increasingly — conduct disruptive or destructive attacks against adversaries' economic and operational capabilities.}
\paragraph{\textbf{Hacktivists} — groups motivated by ideological, political, or social objectives who use cyber attacks as a form of protest or disruption. Hacktivism is generally less sophisticated than financially motivated or nation-state attacks, but the reputational and operational disruption it can cause — particularly through distributed denial of service attacks and website defacement — can be significant.}
\paragraph{\textbf{Insider threats} — employees, contractors, and other individuals with legitimate access to the organisation's systems who misuse that access, whether deliberately or inadvertently. Insider threats are particularly challenging to manage because they operate within the perimeter of conventional security controls. Malicious insiders — those deliberately stealing data, sabotaging systems, or facilitating external attacks — represent a serious risk, particularly where individuals have privileged access to sensitive systems or data. Inadvertent insiders — those whose mistakes, carelessness, or susceptibility to social engineering create security vulnerabilities — are statistically more prevalent and, in aggregate, equally significant.}
\paragraph{\textbf{Opportunistic attackers} — individuals or groups that exploit widely known vulnerabilities, conduct automated scanning for unpatched systems, or use readily available attack tools to target any vulnerable organisation — without specifically targeting the victim. The prevalence of opportunistic attack means that any organisation with internet-connected systems and unaddressed vulnerabilities is a potential target — regardless of its size, its profile, or the apparent value of its data.}

\subsection{Attack Vectors — How Threats Are Realised}
\paragraph{Understanding the methods by which threat actors attempt to exploit vulnerabilities — the attack vectors — is equally important for effective risk assessment. The most significant attack vectors include:}
\paragraph{\textbf{Phishing and social engineering} — attacks that manipulate people into taking actions that compromise security — clicking malicious links, entering credentials on fake websites, transferring funds to fraudulent accounts, or providing access to systems or data. Phishing is the most prevalent initial access vector for the majority of significant cyber incidents — including ransomware attacks, business email compromise fraud, and data breaches. The effectiveness of phishing attacks — despite the significant investment many organisations have made in staff awareness training — reflects both the sophistication of modern phishing techniques and the fundamental challenge of consistently good human decision-making under cognitive load.}
\paragraph{\textbf{Exploitation of software vulnerabilities} — attacks that exploit known or unknown weaknesses in software — operating systems, applications, network devices, and web services — to gain unauthorised access or to execute malicious code. The window of exposure between the disclosure of a software vulnerability and the patching of that vulnerability across an organisation's systems is a critical period of elevated risk — and organisations with poor patch management processes face persistent and significant exposure.}
\paragraph{\textbf{Credential theft and abuse} — attacks that obtain valid user credentials — through phishing, through data breaches that expose passwords used across multiple services, or through brute-force attacks — and use them to authenticate as legitimate users. Multi-factor authentication is one of the most effective controls against credential-based attacks, but its adoption remains incomplete in many organisations.}
\paragraph{\textbf{Supply chain attacks} — attacks that target the organisation indirectly by compromising its suppliers, software vendors, or managed service providers. The SolarWinds attack — in which a widely used network monitoring software update was compromised by a nation-state actor, providing access to thousands of organisations globally — demonstrated the potential scale and sophistication of supply chain attacks. For organisations that depend on third-party software and services — which in practice means every organisation — supply chain security is a significant and challenging risk management dimension.}
\paragraph{\textbf{Ransomware} — as noted above, ransomware has become the most operationally and financially significant cyber threat for most organisations. Modern ransomware attacks are typically multi-stage — beginning with initial access through phishing or vulnerability exploitation, followed by lateral movement across the network to maximise the scope of encryption, followed by data exfiltration before encryption to create additional leverage, and culminating in the deployment of encryption across the organisation's systems. The combination of operational disruption — systems unavailable whilst encrypted — and data breach consequences — personal or sensitive data exfiltrated and threatened for release — makes ransomware one of the most serious risk events an organisation can experience.}

%---------------------------- SECTION ----------------------------
\section{The Data Protection Dimension}
\paragraph{Cybersecurity risk and data protection risk are deeply intertwined — and for compliance professionals, understanding the connection between them is essential. A cybersecurity incident that involves the unauthorised access to, or exfiltration of, personal data is simultaneously a cyber incident and a data breach — triggering both cybersecurity incident response obligations and data protection regulatory obligations.}

\subsection{The GDPR and Data Security}
\paragraph{The \textbf{General Data Protection Regulation} — and the \textbf{UK GDPR} that forms part of UK law through the Data Protection Act 2018 — imposes specific obligations around the security of personal data. Article 32 requires controllers and processors to implement appropriate technical and organisational security measures — taking into account the state of the art, the costs of implementation, and the nature, scope, context, and purposes of processing as well as the risks of varying likelihood and severity to the rights and freedoms of natural persons.}
\paragraph{This is, at its core, a risk assessment obligation. The "appropriate" security measures required by Article 32 are defined by reference to the risks — meaning that organisations need to assess the risks to personal data, and implement security measures that are proportionate to those risks. An organisation that implements generic security controls without conducting a genuine assessment of the specific risks to its personal data processing is unlikely to satisfy this requirement.}
\paragraph{The GDPR specifically mentions several security measures as examples of appropriate controls — pseudonymisation and encryption of personal data, the ability to ensure ongoing confidentiality, integrity, availability, and resilience of processing systems and services, the ability to restore availability and access to personal data in a timely manner in the event of a physical or technical incident, and a process for regularly testing, assessing, and evaluating the effectiveness of technical and organisational measures.}
\paragraph{For compliance professionals, the GDPR's security requirements are particularly significant because they are connected to significant enforcement risk. The ICO has issued substantial fines for data security failures — including for inadequate security measures, failure to implement multi-factor authentication, and failure to patch known vulnerabilities. The enforcement record in this area demonstrates that the ICO treats data security risk assessment as a genuine expectation rather than a best-practice recommendation.}

\subsection{Data Protection Impact Assessments — DPIA}
\paragraph{The \textbf{Data Protection Impact Assessment (DPIA)} — a formal requirement under Article 35 of the GDPR — is the most directly relevant risk assessment obligation in the data protection framework. A DPIA is required when a type of processing is likely to result in a high risk to the rights and freedoms of individuals — with specific categories of processing identified as presumptively requiring a DPIA, including:.}
\begin{itemize}
    \item{Systematic and extensive evaluation of personal aspects, including profiling.}
    \item{Large-scale processing of special category data.}
    \item{Systematic monitoring of publicly accessible areas.}
    \item{Processing involving new technologies.}
    \item{Processing that involves matching or combining datasets.}
    \item{Processing of data concerning vulnerable individuals.}
\end{itemize}
\paragraph{A DPIA is, in substance, a risk assessment — it requires the controller to describe the processing, assess the necessity and proportionality of the processing, assess the risks to the rights and freedoms of individuals, and identify the measures to address those risks. Where the DPIA identifies a high residual risk that cannot be mitigated by the controller, the controller must consult the ICO before proceeding with the processing.}
\paragraph{For compliance professionals, the DPIA obligation creates a direct and recurring risk assessment responsibility — one that is triggered by new processing activities, new technologies, significant changes to existing processing, and periodic review of existing high-risk processing. Building a systematic DPIA process — with clear triggers, defined methodology, appropriate governance, and adequate documentation — is a core data protection compliance obligation.}

\subsection{Personal Data Breach — Notification and Response}
\paragraph{The GDPR imposes specific obligations in relation to personal data breaches — the unauthorised access, loss, destruction, or disclosure of personal data. Controllers are required to:}
\begin{itemize}
    \item{\textbf{Notify the ICO within 72 hours} of becoming aware of a breach that is likely to result in a risk to the rights and freedoms of individuals.}
    \item{\textbf{Notify affected individuals without undue delay} where the breach is likely to result in a high risk to their rights and freedoms.}
\end{itemize}
\paragraph{These notification obligations create a direct and time-sensitive compliance requirement in the event of a cyber incident involving personal data. The 72-hour notification window is short — particularly given the complexity of investigating a significant cyber incident — and organisations that do not have robust incident detection and response capabilities, and clear internal escalation processes, are likely to find themselves in breach of the notification requirement as a consequence of being in breach of their security obligations.}
\paragraph{The ICO's enforcement approach makes clear that late notification, incomplete notification, and failure to notify are treated seriously — and that the circumstances of the breach, the response of the controller, and the quality of the organisation's pre-breach security arrangements are all factors in the ICO's assessment of the appropriate regulatory response.}

%---------------------------- SECTION ----------------------------
\section{Cybersecurity Risk Assessment — The Practical Framework}
\paragraph{With the threat landscape and the regulatory context established, what does a practical, comprehensive cybersecurity risk assessment framework look like?}

\subsection{Identifying Cybersecurity Assets and Dependencies}
\paragraph{The starting point for cybersecurity risk assessment is understanding what needs to be protected — the assets that, if compromised, would result in harm. In a cybersecurity context, the relevant assets include:}
\paragraph{\textbf{Data} — particularly personal data subject to GDPR obligations, commercially sensitive data, intellectual property, financial data, and operational data. Understanding where sensitive data is held, how it flows through the organisation's systems and processes, and who has access to it is the foundation of data-centric security risk assessment.}
\paragraph{\textbf{Systems and applications} — the technology infrastructure through which the organisation's operations are conducted. Understanding the criticality of individual systems — which systems, if unavailable, would cause the most significant operational disruption — is essential for prioritising security investment and business continuity planning.}
\paragraph{\textbf{Operational technology} — for organisations with industrial control systems, building management systems, or other operational technology environments, the security of those systems is a distinct and increasingly important risk domain. The convergence of IT and OT environments has expanded the attack surface for many industrial and infrastructure organisations in ways that their security frameworks may not yet adequately address}
\paragraph{\textbf{Third-party dependencies} — as discussed above, supply chain and third-party security is a significant dimension of the cyber risk landscape. Understanding which third parties have access to the organisation's systems or data, and what the security implications of that access are, is an essential component of a complete cybersecurity risk assessment.}

\subsection{The Cybersecurity Risk Assessment Process}
\paragraph{Applying the five-step process from Part III to the cybersecurity context:}
\paragraph{\textbf{Identifying hazards and risk sources} — in cybersecurity terms, this means identifying the threats relevant to the organisation and the vulnerabilities that those threats could exploit. Threat identification draws on threat intelligence — information about the tactics, techniques, and procedures of threat actors relevant to the organisation's sector and profile — as well as on vulnerability assessment — systematic identification of weaknesses in the organisation's systems, processes, and people.}
\paragraph{Useful tools include:}
\begin{itemize}
    \bitem{Vulnerability scanning and penetration testing}{technical assessments that identify exploitable weaknesses in systems and applications.}
    \bitem{Threat intelligence services}{commercial and open-source intelligence on the threat landscape relevant to the organisation's sector and geography.}
    \bitem{Red team exercises}{simulated attack exercises that test the organisation's detection and response capabilities against realistic adversary behaviour.}
    \bitem{Phishing simulations}{controlled tests of staff susceptibility to phishing attacks that identify training needs and measure the effectiveness of awareness programmes.}
\end{itemize}
\paragraph{\textbf{Determining who and what is affected} — understanding which assets, systems, and data are exposed to each identified threat, and the nature of the harm that could result. This maps directly to the exposure analysis discussed in Chapter~\ref{chap:Step 2 — Determining Who or What Is Affected} — identifying not just that a risk exists but who and what bears the consequences if it materialises.}
\paragraph{\textbf{Evaluating the risks} — assessing the likelihood that each identified threat will successfully exploit the relevant vulnerability, and the severity of the consequences if it does. Cybersecurity risk evaluation draws on threat intelligence, vulnerability severity ratings, asset criticality assessments, and the effectiveness of existing controls — producing a prioritised picture of the risks that most warrant attention and investment.}
\paragraph{\textbf{Deciding on controls} — selecting and implementing security controls that address the most significant identified risks. The \textbf{NIST Cybersecurity Framework} — introduced in Chapter~\ref{chap:Risk Management Frameworks} — provides a widely adopted structure for cybersecurity controls organised around five functions: \textbf{Identify, Protect, Detect, Respond, and Recover}. The Hierarchy of Controls principle applies directly — technical controls embedded in systems and architecture are more reliable than administrative controls that depend on user behaviour}
\paragraph{\textbf{Recording, reviewing, and monitoring} — maintaining current documentation of the cybersecurity risk assessment, monitoring the threat landscape for emerging risks, testing controls regularly, and ensuring that the assessment is updated in response to changes in the threat environment, the organisation's technology estate, or its business operations.}

\subsection{The NIST Cybersecurity Framework in Practice}
\paragraph{The NIST Cybersecurity Framework (CSF) — widely adopted as a reference framework for cybersecurity risk management across sectors and jurisdictions — organises cybersecurity capabilities around its five core functions:}
\paragraph{\textbf{Identify} — developing an organisational understanding of cybersecurity risk to systems, people, assets, data, and capabilities. This encompasses asset management, business environment understanding, governance, risk assessment, and risk management strategy. The Identify function is the foundation — without a clear understanding of what exists, what it does, and how it connects, the remaining functions cannot be adequately scoped.}
\paragraph{\textbf{Protect} — implementing appropriate safeguards to ensure delivery of critical services and to limit the impact of a cybersecurity event. This encompasses access control, awareness and training, data security, information protection processes and procedures, maintenance, and protective technology. The Protect function addresses the prevention dimension of cybersecurity — the controls designed to prevent threats from successfully exploiting vulnerabilities.}
\paragraph{\textbf{Detect} — implementing appropriate activities to identify the occurrence of a cybersecurity event. This encompasses anomaly and event detection, security continuous monitoring, and detection processes. The Detect function is critical because prevention controls are imperfect — the ability to detect a security incident quickly, before an attacker has had time to cause maximum damage, is one of the most important determinants of incident outcomes.}
\paragraph{\textbf{Respond} — implementing appropriate activities to take action regarding a detected cybersecurity incident. This encompasses response planning, communications, analysis, mitigation, and improvements. The quality of the incident response capability — the speed and effectiveness with which the organisation contains and remediates a security incident — directly determines the extent of the harm caused.}
\paragraph{\textbf{Recover} — implementing appropriate activities to maintain plans for resilience and to restore any capabilities or services that were impaired due to a cybersecurity incident. This encompasses recovery planning, improvements, and communications. The Recover function addresses the business continuity dimension of cybersecurity — the ability to restore normal operations after an incident and to learn from the experience to improve future resilience.}

\subsection{ISO 27001 — The Certifiable Standard}
\paragraph{\textbf{ISO 27001} — introduced in Chapter~\ref{chap:Risk Management Frameworks} — provides the international standard for information security management systems, and is the most widely adopted certifiable framework for cybersecurity governance. Certification to ISO 27001 provides external validation that an organisation's information security management meets the standard's requirements — a credential that carries increasing weight with customers, partners, regulators, and insurers.}
\paragraph{The ISO 27001 standard requires organisations to:}
\begin{itemize}
    \item{Establish the context and scope of the information security management system.}
    \item{Conduct a systematic information security risk assessment — identifying risks to the confidentiality, integrity, and availability of information assets.}
    \item{Select and implement controls to address identified risks — drawn from the standard's Annex A control set, which provides a comprehensive library of security controls across 14 domains.}
    \item{Monitor and measure the performance of the information security management system.}
    \item{Continuously improve the system in response to changing risks and lessons learned.}
\end{itemize}
\paragraph{For compliance professionals, ISO 27001 certification provides a robust framework for cybersecurity risk assessment and management — and its alignment with the GDPR's requirements around appropriate technical and organisational security measures makes it a particularly relevant credential for organisations with significant personal data processing activities.}

%---------------------------- SECTION ----------------------------
\section{Sector-Specific Cybersecurity Requirements}
\paragraph{Beyond the general frameworks of GDPR, NIST, and ISO 27001, a number of sector-specific regulatory requirements around cybersecurity risk assessment and management are particularly relevant to compliance professionals in regulated industries.}

\subsection{Financial Services — DORA}
\paragraph{The \textbf{Digital Operational Resilience Act (DORA)} — effective from January 2025 across the EU financial services sector — represents the most comprehensive and most technically detailed cybersecurity regulatory framework yet introduced for financial services. DORA imposes specific requirements on banks, insurers, investment firms, payment institutions, and a wide range of other financial entities around:}
\begin{itemize}
    \bitem{ICT risk management}{a comprehensive framework for identifying, assessing, and managing information and communication technology risks.}
    \bitem{ICT incident classification and reporting}{detailed requirements for classifying, managing, and reporting significant ICT incidents to regulators.}
    \bitem{Digital operational resilience testing}{requirements for regular testing of ICT systems and security, including advanced threat-led penetration testing for significant firms.}
    \bitem{Third-party ICT risk management}{specific requirements for managing the risks associated with third-party technology providers, including contractual requirements and concentration risk monitoring.}
    \bitem{Information sharing}{arrangements for sharing cyber threat intelligence across the financial sector.}
\end{itemize}
\paragraph{For compliance professionals in EU financial services firms — and for non-EU firms with significant EU operations — DORA represents a substantial new compliance obligation that requires a comprehensive review of existing ICT risk management and cybersecurity frameworks.}
\paragraph{In the UK, the FCA and PRA have published operational resilience requirements — including the expectation that firms identify their important business services, set impact tolerances, and test their ability to remain within those tolerances during disruption — that similarly address the cybersecurity and ICT resilience dimension of financial services regulation.}

\subsection{Healthcare}
\paragraph{Healthcare organisations carry a distinctive cybersecurity risk profile — combining the operational criticality of clinical systems (where unavailability can directly threaten patient safety) with the sensitivity of the personal data held (special category health data subject to enhanced GDPR protections) and the attractiveness of that data to threat actors (healthcare data commands premium prices on criminal markets).}
\paragraph{The \textbf{NHS and healthcare sector} has been a prominent target for ransomware attacks — with the 2017 \textbf{\textit{WannaCry}} attack on the NHS demonstrating the devastating operational consequences of a large-scale ransomware incident in a healthcare context. The regulatory framework for NHS cybersecurity — including requirements under the \textbf{Network and Information Systems (NIS) Regulations} and NHS-specific security standards — reflects the critical nature of healthcare infrastructure.}
\paragraph{For compliance professionals in healthcare organisations, the cybersecurity risk assessment framework needs to explicitly address the patient safety dimension of cyber risk — the risk that a cyber incident affecting clinical systems could result in harm to patients — as well as the data protection and regulatory compliance dimensions.}

\subsection{Critical National Infrastructure}
\paragraph{Operators of \textbf{critical national infrastructure} — energy, water, transport, telecoms, and other essential services — are subject to the \textbf{NIS Regulations} in the UK (implementing the EU's Network and Information Systems Directive), which impose security requirements and incident reporting obligations on operators of essential services. The NIS2 Directive — which significantly strengthened and expanded the scope of the original NIS Directive across the EU — reflects the growing regulatory recognition of the systemic risks posed by cyber attacks on critical infrastructure.}

%---------------------------- SECTION ----------------------------
\section{Third-Party Cyber Risk — The Extended Attack Surface}
\paragraph{As noted in the discussion of attack vectors, third-party and supply chain cyber risk is one of the most significant and most challenging dimensions of the cybersecurity risk landscape. For compliance professionals, third-party cyber risk sits at the intersection of cybersecurity risk assessment and third-party risk management — a topic explored in depth in Chapter~\ref{chap:Third-Party and Supply Chain Risk Assessment}.}
\paragraph{The key principles for managing third-party cyber risk include:}
\paragraph{\textbf{Due diligence before onboarding} — assessing the cybersecurity posture of third parties before granting them access to systems or data. This should include a review of the third party's security policies and practices, any relevant certifications, their incident history, and their approach to security in their own supply chain.}
\paragraph{\textbf{Contractual security requirements} — embedding specific cybersecurity requirements in contracts with third parties — including minimum security standards, audit rights, incident notification obligations, and requirements to maintain security certifications.}
\paragraph{\textbf{Ongoing monitoring} — maintaining visibility of the cybersecurity posture of key third parties on an ongoing basis — not just at the point of initial due diligence. This might include periodic reassessment, monitoring of publicly available threat intelligence relating to key suppliers, and tracking of security incidents affecting the third-party ecosystem.}
\paragraph{\textbf{Concentration risk management} — understanding the degree to which the organisation is dependent on specific third-party technology providers — cloud platforms, software vendors, managed security service providers — and managing the concentration risk that arises from over-dependence on a small number of critical suppliers.}

%---------------------------- SECTION ----------------------------
\section{Incident Response — When Controls Fail}
\paragraph{However effective the organisation's preventive controls, the reality is that cyber incidents will occur. The quality of the organisation's incident response capability — its ability to detect incidents quickly, contain them effectively, communicate appropriately, and recover operations — is as important a risk management investment as any preventive control.}
\paragraph{An effective cybersecurity incident response plan addresses:}
\paragraph{\textbf{Detection and initial assessment} — how the organisation identifies that an incident has occurred, assesses its initial scope and severity, and determines the appropriate response level. The faster an incident is detected and the more accurately it is initially assessed, the better the outcomes are likely to be.}
\paragraph{\textbf{Containment} — the immediate actions taken to limit the spread of an incident — isolating affected systems, revoking compromised credentials, blocking malicious network traffic. Effective containment limits the scope of damage and preserves the organisation's ability to investigate and recover.}
\paragraph{\textbf{Eradication} — removing the threat from the organisation's environment — eliminating malware, closing exploited vulnerabilities, and ensuring that the attacker no longer has access to the organisation's systems.}
\paragraph{\textbf{Recovery} — restoring affected systems and data to normal operation. The organisation's backup and recovery capabilities — the quality and currency of backups, the tested ability to restore from those backups, and the time required to do so — directly determine how quickly normal operations can be resumed after a ransomware attack or other destructive incident.}
\paragraph{\textbf{Communication} — managing internal and external communications during and after an incident. Internal communications ensure that the right people have the information they need to make decisions and take action. External communications — to regulators, to affected individuals, to customers, to the media — need to be timely, accurate, and carefully managed. The 72-hour GDPR notification requirement imposes a specific and non-negotiable external communication obligation where personal data is involved.}
\paragraph{\textbf{Post-incident review} — learning from the incident. Every significant cyber incident should prompt a structured review of what happened, why existing controls failed, what the response did well and where it fell short, and what improvements are needed. The outputs of post-incident review should feed directly back into the cybersecurity risk assessment framework — updating threat assessments, control evaluations, and response plans.}

%---------------------------- SECTION ----------------------------
\section{The Board and Cyber Risk — Governance in Practice}
\paragraph{For compliance professionals, one of the most important dimensions of cybersecurity risk management is ensuring that it receives adequate board-level governance. Cyber risk has historically been underrepresented at board level — either because it was perceived as too technical for non-specialist board members, or because it was delegated entirely to IT and security functions without adequate governance oversight.}
\paragraph{Regulators are increasingly explicit that this is not acceptable. The FCA, the PRA, DORA, and the SEC's climate disclosure rules all reflect a growing regulatory expectation that boards engage genuinely with cyber and technology risk — that they maintain adequate visibility of the organisation's cyber risk profile, that they ensure adequate resources are invested in cybersecurity, and that they are appropriately involved in significant cyber risk decisions.}
\paragraph{For compliance professionals supporting board cyber governance, the key contributions are:}
\paragraph{\textbf{Making cyber risk accessible} — translating technical cybersecurity concepts into business risk terms that board members without technical backgrounds can engage with. The board does not need to understand the technical details of a particular vulnerability — it needs to understand what business harm could result if it were exploited, and what investment is required to address it.}
\paragraph{\textbf{Providing honest assurance} — giving the board an accurate picture of the organisation's cyber risk profile and the effectiveness of its controls — including honest acknowledgement of gaps, weaknesses, and areas of uncertainty. A board that is consistently told that cyber risk is under control will not engage with it effectively — a board that is given an honest account of the risk landscape, including its most challenging dimensions, is better placed to fulfil its governance responsibilities.}
\paragraph{\textbf{Connecting cyber risk to the enterprise risk framework} — ensuring that cyber risk is assessed and reported within the same governance framework as the organisation's other material risks — not in a separate IT committee report that is disconnected from the main board risk agenda.}

%---------------------------- SECTION ----------------------------
\newpage
\section{Chapter Summary}
In this chapter we learned about the following key points:

\begin{table}[H]
  \centering
  \renewcommand{\arraystretch}{1.5}
  \begin{tabularx}{\textwidth}{ | >{\raggedright\arraybackslash}p{0.35\textwidth} 
								| >{\raggedright\arraybackslash}p{0.35\textwidth} 
                                | >{\raggedright\arraybackslash}X | }
    \hline
    \textbf{Threat Category} & \textbf{Primary Mechanism} & \textbf{Key Controls} \\
    \hline
    \textbf{Financially motivated criminals} & Ransomware; business email compromise; fraud & MFA; patching; backups; staff awareness\\
    \hline
    \textbf{Nation-state actors} & Advanced persistent threat; supply chain compromise & Threat intelligence; network segmentation; privileged access management\\
    \hline
    \textbf{Insider threats} & Data theft; sabotage; inadvertent exposure & Access controls; monitoring; culture; background screening\\
    \hline
    \textbf{Opportunistic attackers} & Vulnerability exploitation; credential stuffing & Patching; vulnerability management; MFA\\
    \hline
    \textbf{Supply chain attacks} & Third-party compromise; software supply chain & Third-party due diligence; software assurance; contractual requirements\\
    \hline
  \end{tabularx}
  \caption{Cybersecurity Threat Categories}
\end{table}

\begin{table}[H]
  \centering
  \renewcommand{\arraystretch}{1.5}
  \begin{tabularx}{\textwidth}{ | >{\raggedright\arraybackslash}p{0.35\textwidth} 
								| >{\raggedright\arraybackslash}p{0.35\textwidth} 
                                | >{\raggedright\arraybackslash}X | }
    \hline
    \textbf{Function} & \textbf{Purpose} & \textbf{Examples} \\
    \hline
    \textbf{Identify} & Understand the risk landscape & Asset inventory; risk assessment; threat intelligence\\
    \hline
    \textbf{Protect} & Prevent successful attacks & Access control; encryption; patching; training\\
    \hline
    \textbf{Detect} & Identify incidents quickly & SIEM; anomaly detection; monitoring\\
    \hline
    \textbf{Respond} & Contain and manage incidents & Incident response plan; crisis communications; regulatory notification\\
    \hline
    \textbf{Recover} & Restore operations & Backup and recovery; business continuity; post-incident review\\
    \hline
  \end{tabularx}
  \caption{The NIST Cybersecurity Framework}
\end{table}

\begin{table}[H]
  \centering
  \renewcommand{\arraystretch}{1.5}
  \begin{tabularx}{\textwidth}{ | >{\raggedright\arraybackslash}p{0.30\textwidth} 
								| >{\raggedright\arraybackslash}p{0.25\textwidth} 
                                | >{\raggedright\arraybackslash}X | }
    \hline
    \textbf{Obligation} & \textbf{Source} & \textbf{Key Requirement} \\
    \hline
    \textbf{Data security} & GDPR Article 32 & Appropriate technical and organisational security measures\\
    \hline
    \textbf{DPIA} & GDPR Article 35 & Risk assessment for high-risk processing activities\\
    \hline
    \textbf{Breach notification} & GDPR Articles 33-34 & 72-hour notification to ICO; notification to individuals for high-risk breaches\\
    \hline
    \textbf{ICT risk management} & DORA & Comprehensive ICT risk framework; testing; third-party management\\
    \hline
    \textbf{Operational resilience} & FCA/PRA & Important business services; impact tolerances; resilience testing\\
    \hline
    \textbf{NIS Regulations} & NIS/NIS2 & Security requirements and incident reporting for essential services\\
    \hline
  \end{tabularx}
  \caption{Key Regulatory Obligations}
\end{table}

\paragraph{Key takeaways:}
\begin{itemize}
  \item{Cybersecurity risk assessment is not a technical exercise for IT departments — it is a core compliance and governance obligation requiring genuine board engagement and second-line oversight.}
  \item{The threat landscape is continuously evolving — phishing, ransomware, supply chain attacks, and insider threats are the most prevalent current threats but the picture will change, and the risk assessment framework must change with it.}
  \item{Data protection and cybersecurity are deeply interconnected — a cyber incident involving personal data is simultaneously a security incident and a data breach, triggering both cybersecurity and GDPR obligations.}
  \item{The GDPR's Article 32 security requirement and the DPIA obligation are direct, enforceable risk assessment obligations — not best-practice recommendations.}
  \item{DORA represents a significant new regulatory framework for financial services that raises the bar for ICT risk management, testing, and third-party oversight substantially.}
  \item{Third-party cyber risk is among the most challenging dimensions of the cybersecurity risk landscape — due diligence, contractual requirements, and ongoing monitoring are all essential.}
  \item{Incident response capability is as important as preventive controls — the ability to detect, contain, recover, and learn from incidents determines outcomes when prevention fails.}
  \item{The 72-hour GDPR breach notification window is short and non-negotiable — organisations without robust detection and escalation processes will struggle to comply.}
\end{itemize}

\barquote{In Chapter~\ref{chap:Third-Party and Supply Chain Risk Assessment}, we focus specifically on third-party and supply chain risk assessment — exploring the full scope of the obligation, the analytical frameworks available, and how to build a proportionate, effective third-party risk management programme that satisfies both regulatory expectations and genuine management needs.}